\appendix
\setcounter{tocdepth}{1}
\setcounter{secnumdepth}{0}

\section{Introduzione}
Il presente documento costituisce il \textbf{Glossario} ufficiale del gruppo \textbf{Git Happens} per il progetto \textbf{Sonar} (capitolato C8, proponente \textit{Halue S.r.l. SB}). 
L'obiettivo di questa sezione è fornire definizioni chiare, univoche e rigorose di tutti i termini tecnici, gli acronimi, le abbreviazioni e i concetti chiave che potrebbero risultare ambigui o complessi all'interno della documentazione di progetto. 
All'interno dei documenti redatti dal gruppo, ciascun termine incluso nel glossario è contrassegnato da una lettera pedice ``G'' (ad esempio, \term{Termine}) per indicare immediatamente al lettore la presenza di una definizione formale e condivisa, facilitando la comprensione e standardizzando il linguaggio utilizzato tra tutti gli stakeholder coinvolti.
\section{A}
\glossentry{Analisi dei requisiti}{}
L'attività di analisi dei requisiti consiste nella comprensione e formalizzazione delle esigenze espresse dall'azienda proponente Halue S.r.l. SB per il progetto. I requisiti sono classificati in obbligatori, desiderabili e opzionali, distinguendo tra aspetti funzionali, di qualità e di vincolo. L'intero processo è tracciato mediante matrici di copertura che collegano le richieste con i casi d'uso in \term{Unified Modeling Language} e con i relativi test, garantendo coerenza durante tutto lo \term{Sviluppo}.
\glossentry{Attività}{}
Le \term{Attività} di progetto rappresentano l'insieme delle unità di lavoro operative necessarie per raggiungere gli \term{Obiettivo} stabiliti. Gestite tramite l'\term{Issue Tracking System} di \term{Jira} e coordinate dal \term{Team} in base ai rispettivi \term{Ruolo}, esse si suddividono in compiti di gestione, analisi e progettazione, programmazione e implementazione (tramite \term{Proof of Concept} e \term{Minimum Viable Product}), e verifica e controllo mediante \term{Pull Request}. L'esecuzione è organizzata in iterazioni bisettimanali basate su \term{Task} definiti.
\section{B}
\glossentry{Budget}{}
Il \term{Budget} rappresenta la stima preventiva delle risorse finanziarie necessarie per il completamento del progetto o di specifiche \term{Task}, allocate e monitorate dal \term{Team} in coerenza con la \term{Technology Baseline} e gli \term{Obiettivo} stabiliti.
\section{C}
\glossentry{Continuous Integration}{CI}
La \term{Continuous Integration} è una pratica di sviluppo software che prevede l'integrazione frequente e automatica del codice sorgente in un repository condiviso, accompagnata da test automatizzati per rilevare tempestivamente eventuali conflitti o errori.
\glossentry{Cost Performance Index}{CPI}
Il \term{Cost Performance Index} è una metrica di project management utilizzata per valutare l'efficienza dei costi di un progetto, confrontando il valore del lavoro effettivamente completato con i costi effettivi sostenuti.
\glossentry{Costo}{costi}
Il \term{Costo} rappresenta il consumo di risorse finanziarie, temporali o umane impiegate per lo svolgimento delle \term{Task} e il raggiungimento degli \term{Obiettivo} di progetto, monitorato costantemente dal \term{Team}.
\glossentry{Customer Acceptance}{}
La \term{Customer Acceptance} è la fase in cui il prodotto software completato viene presentato al committente per verificarne la conformità ai requisiti concordati e ottenere l'approvazione formale.
\section{D}
\glossentry{Dashboard}{}
La \term{Dashboard} è un'interfaccia grafica centralizzata che aggrega, elabora e visualizza in tempo reale le metriche chiave di progetto, gli indicatori di prestazione e lo stato di avanzamento delle \term{Task}. Questo strumento, introdotto a supporto del monitoraggio operativo del \term{Team}, è stato mutuato direttamente dalle specifiche e dalle esigenze descritte nei documenti di progetto forniti dal proponente.
\glossentry{Data Transfer Object}{DTO}
Il \term{Data Transfer Object} (DTO) è un pattern architetturale utilizzato per trasferire dati tra sottosistemi o componenti software differenti. Il suo impiego evita di esporre direttamente i modelli di dominio interni, riducendo l'accoppiamento e ottimizzando le performance di rete e serializzazione, in conformità con i requisiti architetturali definiti per il \term{Sviluppo}.
\glossentry{Demo}{}
La \term{Demo} (dimostrazione) è una presentazione pratica e funzionale del prodotto software (o di un suo prototipo incrementale) rivolta al committente e agli stakeholder, finalizzata a mostrare le funzionalità implementate e a raccogliere feedback tempestivi durante il ciclo di vita.
\glossentry{DevOps}{}
Il termine \term{DevOps} indica l'insieme di pratiche, strumenti e filosofie culturali che uniscono lo \term{Sviluppo} software (Dev) e le operazioni IT (Ops). L'obiettivo è automatizzare e accelerare i processi di integrazione, testing e rilascio, garantendo elevati standard di qualità e stabilità del sistema in linea con i processi adottati dal gruppo.
\glossentry{Documento}{documenti}
Il \term{Documento} rappresenta qualsiasi artefatto informativo formale redatto dal \term{Team} (quali specifiche, verbali, piani di progetto o manuali) strutturato secondo le norme redazionali stabilite e tracciato all'interno del sistema di gestione documentale.
\section{E}
\glossentry{Earned Value}{EV}
L'\term{Earned Value} (valore guadagnato) è una metrica fondamentale del project management che esprime in termini monetari o di risorse il valore del lavoro effettivamente completato rispetto alla pianificazione iniziale, fornendo una misura oggettiva dell'avanzamento reale del progetto.
\glossentry{Earned Value Management}{EVM}
L'\term{Earned Value Management} (EVM) è una metodologia di gestione e controllo integrato dei progetti che unisce ambito, tempi e costi. Attraverso il monitoraggio di indicatori chiave, consente di misurare le performance e di prevedere con precisione l'andamento futuro e la spesa finale del progetto, in conformità con i processi definiti nei documenti di pianificazione.
\glossentry{Estimate at Completion}{EAC}
L'\term{Estimate at Completion} (EAC) rappresenta la stima aggiornata del costo totale che il progetto avrà una volta completato, calcolata periodicamente a partire dai costi effettivi sostenuti e dalle performance registrate dal \term{Team}.
\glossentry{Estimate to Complete}{ETC}
L'\term{Estimate to Complete} (ETC) indica la stima delle risorse economiche e temporali ancora necessarie per portare a termine tutte le \term{Task} rimanenti del progetto, derivata dalle analisi di bilancio e dallo stato di avanzamento corrente.
\section{F}
\glossentry{Fixture}{fixtures}
Nel contesto dello \term{Sviluppo} software e dei test automatici, una \term{Fixture} rappresenta lo stato iniziale fisso o un insieme di dati predefiniti configurati per garantire che l'ambiente di prova sia stabile, controllato e ripetibile durante l'esecuzione dei test (come quelli effettuati tramite \term{Vitest} o \term{Playwright}). L'uso delle fixture è una pratica standard adottata dal \term{Team} per isolare i test e assicurare la coerenza dei risultati.
\section{G}
\glossentry{Gang of Four}{GoF}
Il termine \term{Gang of Four} si riferisce al gruppo di autori (Erich Gamma, Richard Helm, Ralph Johnson e John Vlissides) che ha codificato per primo i 23 pattern di progettazione software fondamentali (\emph{design patterns}) orientati agli oggetti. Tali pattern costituiscono il riferimento architetturale fondamentale adottato dal \term{Team} durante la progettazione del codice.
\glossentry{Gantt}{}
Il diagramma di \term{Gantt} è uno strumento di pianificazione temporale e visualizzazione grafica della programmazione di progetto. Permette di rappresentare le \term{Task}, la loro durata, le dipendenze reciproche e la distribuzione temporale lungo il ciclo di vita, consentendo al \term{Team} di monitorare costantemente le scadenze e l'avanzamento dei lavori.
\glossentry{Git}{}
\term{Git} è un sistema di controllo di versione distribuito e open source, progettato per gestire con efficacia e velocità qualsiasi tipo di progetto software. Consente a più sviluppatori di lavorare in parallelo sullo stesso codice sorgente, tracciando ogni modifica, gestendo i rami di sviluppo (\emph{branch}) e risolvendo eventuali conflitti d'integrazione.
\glossentry{GitHub}{}
\term{GitHub} è una piattaforma basata sul web per l'hosting di repository \term{Git} e la collaborazione allo \term{Sviluppo} software. Offre strumenti avanzati per la gestione del codice, il tracciamento delle modifiche, la revisione paritaria tramite \emph{Pull Request} e l'automazione dei flussi di lavoro tramite integrazione continua, rappresentando la piattaforma principale utilizzata dal gruppo per la gestione del progetto.
\section{I}
\glossentry{Incontro}{incontri}
L'\term{Incontro} (o riunione) rappresenta un momento di confronto formale o informale tra i membri del \term{Team}, o tra il gruppo e il proponente \term{Halue S.r.l. SB}, finalizzato alla condivisione dello stato di avanzamento, alla discussione di scelte architetturali o organizzative e alla stesura di un apposito \term{Verbale}.
\glossentry{Indice}{indici}
L'\term{Indice} è la struttura di navigazione testuale inserita all'inizio della documentazione di progetto, che elenca in modo gerarchico sezioni, capitoli e contenuti per facilitare la consultazione dei documenti redatti dal gruppo.
\glossentry{Integrated Development Environment}{IDE}
Un \term{Integrated Development Environment} (IDE) è un software applicativo che fornisce un ambiente di \term{Sviluppo} completo, integrando in un'unica interfaccia editor di codice, strumenti di compilazione, debugging e gestione del controllo di versione.
\glossentry{Institute of Electrical and Electronics Engineers}{IEEE}
L'\term{Institute of Electrical and Electronics Engineers} (IEEE) è la più grande organizzazione professionale internazionale senza scopo di lucro dedicata all'avanzamento della tecnologia e alla definizione di standard globali in ambito ingegneristico, informatico ed elettronico, i cui modelli normativi costituiscono il riferimento per i processi di qualità del progetto.
\glossentry{International Organization for Standardization}{ISO}
L'\term{International Organization for Standardization} (ISO) è un'organizzazione non governativa internazionale che sviluppa e pubblica norme tecniche e di gestione a livello globale. Le specifiche ISO (come quelle relative ai cicli di vita del software e alla qualità) rappresentano il quadro normativo di riferimento adottato dal gruppo.
\glossentry{Issue Tracking System}{ITS}
L'\term{Issue Tracking System} è lo strumento software utilizzato per la gestione, l'assegnazione e il monitoraggio delle \term{Task}, dei bug e delle richieste di modifica durante l'intero ciclo di vita del progetto, integrato direttamente nei flussi di lavoro del \term{Team} tramite piattaforme come \term{Jira}.
\section{J}
\glossentry{JavaScript Object Notation}{JSON}
La \term{JavaScript Object Notation} (JSON) è un formato di interscambio di dati leggero, basato su testo e indipendente dal linguaggio di programmazione. Grazie alla sua struttura basata su coppie chiave-valore, viene ampiamente utilizzata nello \term{Sviluppo} software moderno per la comunicazione tra client e server, ad esempio attraverso API REST.
\glossentry{Jira}{}
\term{Jira} è un sistema proprietario di \term{Issue Tracking System} e gestione di progetto sviluppato da Atlassian. Viene utilizzato dal \term{Team} per la pianificazione delle attività, il tracciamento dei \term{Task}, la gestione dei requisiti e il monitoraggio dell'avanzamento dei lavori all'interno delle iterazioni di \term{Scrum} o \term{Kanban}, come previsto dai processi di progetto.
\section{K}
\glossentry{Kanban}{}
Il \term{Kanban} è un metodo di gestione del lavoro visuale e agile, progettato per aiutare i team a identificare, gestire e ottimizzare il flusso di lavoro. Attraverso l'uso di una bacheca digitale o fisica divisa in colonne (che rappresentano i diversi stati delle \term{Task}, come \emph{To Do}, \emph{In Progress} e \emph{Done}), limita il lavoro in corso (\emph{WIP limits}) per massimizzare l'efficienza e ridurre i colli di bottiglia, integrandosi con gli strumenti di project management adottati dal \term{Team}.
\section{L}
\glossentry{LaTeX}{}
\term{LaTeX} è un linguaggio e un sistema di preparazione dei documenti basato su markup, ampiamente utilizzato nell'ambito accademico e professionale per la stesura di testi tecnici e scientifici complessi. Grazie alla separazione tra contenuto e formattazione, garantisce un'elevata qualità tipografica ed è lo strumento principale adottato dal \term{Team} per la redazione della documentazione ufficiale di progetto.
\section{M}
\glossentry{Method Branch Coverage}{}
La \term{Method Branch Coverage} (copertura dei rami per metodo) è una metrica di test strutturale che misura la percentuale di rami decisionali (come i rami \emph{true} e \emph{false} nelle strutture condizionali) che sono stati effettivamente eseguiti durante la suite di test rispetto al totale dei rami presenti in ciascun metodo, garantendo un elevato standard di affidabilità del codice.
\glossentry{Minimum Viable Product}{MVP}
Il \term{Minimum Viable Product} (MVP) rappresenta la versione iniziale di un prodotto software dotata del set minimo di funzionalità indispensabili per soddisfare i requisiti chiave e consentire agli utenti o al committente di testare il valore della soluzione, fungendo da base per i successivi incrementi di \term{Sviluppo}.
\glossentry{Mock}{mocks}
Un \term{Mock} è un oggetto simulato utilizzato nei test automatizzati per replicare il comportamento di componenti reali complessi o esterni (come database o API di terze parti). Permette al \term{Team} di isolare la porzione di codice in esame e verificarne il corretto funzionamento in un ambiente controllato.
\glossentry{Mock Service Worker}{MSW}
Il \term{Mock Service Worker} (MSW) è una libreria moderna utilizzata nello \term{Sviluppo} frontend per intercettare le richieste di rete a livello di livello di servizio (\emph{Service Worker}). Consente di simulare risposte di rete realistiche sia durante lo sviluppo locale che durante l'esecuzione dei test, garantendo stabilità e indipendenza dai sistemi esterni.
\glossentry{Modified Statement Coverage}{}
La \term{Modified Statement Coverage} (copertura delle istruzioni modificata) è una metrica di testing orientata a verificare che le singole istruzioni eseguibili all'interno del codice sorgente siano state coperte in modo adeguato dai test automatici, contribuendo al raggiungimento degli obiettivi di qualità definiti nei piani di verifica.
\section{N}
\glossentry{Next.js}{}
\term{Next.js} è un framework \term{React} open source sviluppato da Vercel che consente la creazione di applicazioni web orientate alla produzione. Offre funzionalità avanzate come il rendering lato server (\emph{SSR}), la generazione di siti statici (\emph{SSG}) e il routing basato su file system, rappresentando la tecnologia principale scelta dal \term{Team} per lo \term{Sviluppo} dell'interfaccia utente del progetto.
\section{O}
\glossentry{Obiettivo}{obiettivi}
L'\term{Obiettivo} rappresenta il traguardo o il risultato atteso che il \term{Team} si impegna a raggiungere, sia in termini di qualità e funzionalità del prodotto software sia rispetto alla pianificazione temporale e al rispetto del \term{Budget} stabiliti nei piani di progetto.
\glossentry{Object-Relational Mapping}{ORM}
L'\term{Object-Relational Mapping} (ORM) è una tecnica di programmazione che consente di convertire i dati tra sistemi di tipo incompatibili nei database relazionali e nei linguaggi di programmazione orientati agli oggetti, facilitando l'interazione con la base di dati tramite oggetti software e accelerando lo \term{Sviluppo}.
\glossentry{Open Authorization}{OAuth, OpenAuthorization}
L'\term{Open Authorization} (OAuth) è un protocollo standard aperto per l'autorizzazione che consente agli utenti di concedere a siti Web o applicazioni di terze parti l'accesso alle proprie informazioni su un altro servizio, senza doverne condividere le credenziali di accesso e garantendo elevati standard di sicurezza.
\glossentry{OpenID Connect}{OIDC}
L'\term{OpenID Connect} (OIDC) è un protocollo di autenticazione basato su standard aperto, costruito come livello sopra il protocollo \term{OpenAuthorization} (OAuth 2.0). Consente alle applicazioni di verificare l'identità dell'utente sulla base dell'autenticazione eseguita da un server di autorizzazione esterno, ricevendo al contempo informazioni di profilo di base tramite un token di sicurezza (\emph{ID Token}).
\section{P}
\glossentry{Plan-Do-Check-Act}{PDCA}
Il ciclo \term{Plan-Do-Check-Act} (PDCA), noto anche come ciclo di Deming, è un modello di gestione iterativo in quattro fasi utilizzato per il controllo e il miglioramento continuo dei processi e della qualità del prodotto software. Adottato dal \term{Team} come riferimento metodologico, consente di pianificare le attività, eseguirle, verificarne i risultati rispetto agli \term{Obiettivo} e attuare azioni correttive mirate.
\glossentry{Planned Value}{PV}
Il \term{Planned Value} (PV) è una metrica chiave dell'\term{Earned Value Management} che rappresenta il valore autorizzato del lavoro pianificato da completare fino a un determinato momento del progetto, fungendo da riferimento (\emph{baseline}) per misurare l'avanzamento temporale ed economico delle \term{Task}.
\glossentry{Playwright}{}
\term{Playwright} è un framework open source sviluppato da Microsoft per l'automazione e il testing \emph{end-to-end} di applicazioni web moderne. Supporta i principali motori di rendering (Chromium, Firefox e WebKit) e viene impiegato dal \term{Team} per verificare il corretto funzionamento dell'interfaccia utente e dei flussi applicativi in scenari reali.
\glossentry{Portable Document Format}{PDF}
Il \term{Portable Document Format} (PDF) è un formato di file standard aperto sviluppato per rappresentare documenti in modo indipendente dall'hardware, dal sistema operativo e dal software applicativo utilizzati per crearli o visualizzarli. È il formato standard adottato dal gruppo per l'esportazione e la consegna ufficiale di ogni \term{Documento} redatto in \term{LaTeX}.
\glossentry{Preventivo}{preventivi}
Il \term{Preventivo} rappresenta la stima formale dei costi, delle risorse umane e dei tempi necessari per la realizzazione del progetto, elaborata dal \term{Team} in fase di pianificazione iniziale e utilizzata come base di riferimento per il monitoraggio del \term{Budget} e la consuntivazione.
\glossentry{Program Evaluation and Review Technique}{PERT}
La \term{Program Evaluation and Review Technique} (PERT) è una tecnica statistica di project management utilizzata per analizzare, rappresentare e stimare la durata delle \term{Task} necessarie al completamento di un progetto, tenendo conto dell'incertezza tramite stime ottimistiche, probabili e pessimistiche.
\glossentry{Proof of Concept}{PoC}
Il \term{Proof of Concept} (PoC) è una realizzazione prototipale parziale o sperimentale sviluppata con l'obiettivo di dimostrare la fattibilità tecnica di una soluzione architetturale o l'idoneità delle tecnologie scelte prima di procedere allo \term{Sviluppo} completo del prodotto.
\glossentry{Pull Request}{PR, pull requests}
Una \term{Pull Request} (PR) è un meccanismo collaborativo offerto da piattaforme come \term{GitHub} che consente a uno sviluppatore di notificare al resto del \term{Team} il completamento di una modifica su un ramo (\emph{branch}) del repository \term{Git}. Prima dell'integrazione nel ramo principale, la richiesta viene sottoposta a revisione paritaria (\emph{code review}) e ai controlli automatici di \term{Continuous Integration} per garantire la qualità del codice.
\section{R}
\glossentry{Refactoring}{}
Il \term{Refactoring} è il processo di ristrutturazione e miglioramento del codice sorgente esistente senza alterarne il comportamento esterno. Viene applicato dal \term{Team} per ridurre la complessità, eliminare la duplicazione e migliorare la leggibilità e la manutenibilità del software durante lo \term{Sviluppo}.
\glossentry{Representational State Transfer}{REST}
Il \term{Representational State Transfer} (REST) è uno stile architetturale per sistemi distribuiti, ampiamente utilizzato nella progettazione di servizi web e API. Basato su operazioni stateless e sull'uso standardizzato dei metodi HTTP per manipolare risorse (spesso scambiate in formato \term{JavaScript Object Notation}), costituisce il modello di riferimento adottato nel progetto per la comunicazione tra client e server.
\glossentry{Requirements and Technology Baseline}{RTB}
La \term{Requirements and Technology Baseline} (RTB) rappresenta una delle principali milestone di progetto, mutuata direttamente dal capitolato e dai documenti di pianificazione. In questa fase viene consolidata la specifica dei requisiti e validata l'architettura tecnologica tramite un \term{Proof of Concept}, verificando l'idoneità degli strumenti scelti prima di procedere allo sviluppo completo del prodotto.
\glossentry{Riunione}{riunioni}
La \term{Riunione} è un momento di incontro strutturato e pianificato tra i membri del \term{Team} (riunione interna) o con il proponente \term{Halue S.r.l. SB} (riunione esterna), finalizzato all'allineamento sullo stato dei lavori, alla risoluzione di criticità e alla presa di decisioni operative, documentate formalmente tramite un \term{Verbale}.
\glossentry{Row Level Security}{RLS}
La \term{Row Level Security} (RLS) è una funzionalità di sicurezza avanzata dei database relazionali che consente di controllare l'accesso alle singole righe di una tabella in base all'identità o al ruolo dell'utente che esegue la query. Questo meccanismo, introdotto per soddisfare i vincoli di riservatezza emersi dalla descrizione del progetto, garantisce che ciascun utente possa visualizzare o modificare esclusivamente i dati di propria competenza.
\glossentry{Ruolo}{ruoli}
Il \term{Ruolo} definisce l'insieme delle responsabilità, delle competenze e dei compiti assegnati a ciascun membro del \term{Team} all'interno del ciclo di vita del progetto (quali Responsabile, Amministratore, Analista, Progettista, Programmatore e Verificatore). La rotazione periodica dei ruoli garantisce un'equa distribuzione del carico di lavoro e l'acquisizione di competenze trasversali da parte di tutto il gruppo.
\section{S}
\glossentry{Scadenza}{scadenze}
La \term{Scadenza} rappresenta il limite temporale prefissato entro il quale una specifica \term{Task}, una milestone o la consegna di un \term{Documento} deve essere completata, monitorata costantemente dal \term{Team} per garantire il rispetto della pianificazione.
\glossentry{Schedule Performance Index}{SPI}
Lo \term{Schedule Performance Index} (SPI) è una metrica dell'\term{Earned Value Management} che misura l'efficienza temporale di un progetto, calcolata come rapporto tra il valore del lavoro effettivamente svolto (\term{Earned Value}) e quello pianificato (\term{Planned Value}). Un valore inferiore a 1 indica un ritardo rispetto alla tabella di marcia.
\glossentry{Schedule Variance}{SV}
La \term{Schedule Variance} (SV) è un indicatore dell'\term{Earned Value Management} che esprime, in termini monetari o di risorse, lo scostamento temporale tra il lavoro realmente completato e quello pianificato fino a una determinata data, consentendo di individuare tempestivamente eventuali ritardi nelle \term{Task}.
\glossentry{Scrum}{}
\term{Scrum} è un framework agile per lo \term{Sviluppo} e la gestione di progetti complessi, basato su iterazioni incrementali a durata fissa dette \emph{Sprint}. Adottato dal \term{Team} come riferimento metodologico insieme a \term{Kanban}, favorisce la collaborazione continua, l'adattabilità ai cambiamenti e il rilascio frequente di valore tramite cerimonie strutturate e \term{Ruolo} ben definiti.
\glossentry{Sezione}{sezioni}
La \term{Sezione} è un'unità logica e strutturale di suddivisione interna di un \term{Documento}, utilizzata per organizzare i contenuti in modo gerarchico e modulare, migliorandone la leggibilità e la tracciabilità.
\glossentry{SOLID}{}
\term{SOLID} è un acronimo che identifica i cinque principi fondamentali della programmazione e della progettazione orientata agli oggetti (\emph{Single Responsibility}, \emph{Open-Closed}, \emph{Liskov Substitution}, \emph{Interface Segregation} e \emph{Dependency Inversion}). Tali principi guidano il \term{Team} nella scrittura di codice pulito, modulare, facilmente estensibile e manutenibile durante tutto lo \term{Sviluppo}.
\glossentry{Stato Avanzamento Lavori}{SAL}
Lo \term{Stato Avanzamento Lavori} (SAL) è un momento di verifica periodica, interno o condiviso con il committente \term{Halue S.r.l. SB}, finalizzato a rendicontare le attività svolte, valutare le metriche di qualità e di costo e riallineare la pianificazione per le iterazioni successive, come previsto dai processi di progetto.
\glossentry{Stub}{stubs}
Uno \term{Stub} è un componente software fittizio utilizzato nei test automatizzati per sostituire una dipendenza reale non ancora implementata o complessa da invocare, restituendo risposte predefinite e controllate al fine di isolare e verificare il comportamento del modulo sotto test.
\glossentry{Supabase}{}
\term{Supabase} è una piattaforma open source \emph{Backend-as-a-Service} (BaaS) basata su PostgreSQL, che fornisce database relazionale, autenticazione, API istantanee e funzionalità di \term{Row Level Security}. Questa tecnologia è stata selezionata nella \term{Requirements and Technology Baseline} per gestire la persistenza e la sicurezza dei dati secondo i vincoli emersi dalla descrizione del progetto.
\glossentry{Sviluppo}{sviluppi}
Lo \term{Sviluppo} rappresenta l'insieme delle attività tecniche e ingegneristiche (analisi, progettazione, codifica e verifica) svolte dal \term{Team} per trasformare i requisiti concordati in un prodotto software funzionante e conforme agli standard qualitativi stabiliti.
\section{T}
\glossentry{Task}{tasks}
Una \term{Task} rappresenta un'unità di lavoro elementare, circoscritta e assegnata a uno o più membri del \term{Team} all'interno della pianificazione di progetto. Ciascuna task è tracciata tramite l'\term{Issue Tracking System} e concorre al completamento degli obiettivi di \term{Sprint} o di intere fasi di \term{Sviluppo}.
\glossentry{Team}{}
Il \term{Team} indica il gruppo di lavoro dei componenti del progetto, organizzato attraverso la definizione di ruoli e responsabilità ben definiti e coordinato per garantire la corretta esecuzione delle attività di analisi, progettazione, codifica, verifica e redazione della documentazione.
\glossentry{Test-Driven Development}{TDD}
Il \term{Test-Driven Development} (TDD) è una pratica di sviluppo software in cui i test automatici vengono scritti prima del codice funzionale vero e proprio. Il processo segue un ciclo iterativo di tre fasi (\emph{red-green-refactoring}): si scrive un test che fallisce, si implementa il codice minimo necessario per superarlo e infine si procede all'\term{Refactoring} del codice per migliorarne la qualità, garantendo un'elevata copertura e robustezza del sistema.
\glossentry{TypeScript}{}
\term{TypeScript} è un linguaggio di programmazione open source sviluppato da Microsoft che estende \term{JavaScript} introducendo la tipizzazione statica opzionale. Viene adottato dal \term{Team} nello \term{Sviluppo} per migliorare la manutenibilità del codice, intercettare errori in fase di compilazione e facilitare la collaborazione su basi di codice complesse.
\section{U}
\glossentry{Unified Modeling Language}{UML}
L'\term{Unified Modeling Language} (UML) è un linguaggio di modellazione visuale standardizzato nell'ingegneria del software, utilizzato per specificare, visualizzare, costruire e documentare gli artefatti di un sistema software. Il \term{Team} lo impiega per redigere diagrammi strutturali e comportamentali (come i diagrammi dei casi d'uso, delle classi e di sequenza) utili a chiarire l'architettura e le interazioni del prodotto prima e durante lo \term{Sviluppo}.
\glossentry{Utente}{utenti}
L'\term{Utente} rappresenta qualsiasi soggetto fisico o entità esterna che interagisce direttamente con il sistema software per usufruire delle sue funzionalità. A seconda del contesto e delle politiche di accesso (come quelle gestite tramite \term{Row Level Security}), gli utenti possono assumere profili e privilegi differenziati in base ai requisiti emersi in fase di analisi.
\section{V}
\glossentry{Verbale}{verbali}
Il \term{Verbale} è il documento ufficiale redatto al termine di ogni \term{Incontro} (sia interno al \term{Team} che esterno con il proponente \term{Halue S.r.l. SB}), contenente la sintesi delle tematiche discusse, le decisioni prese, le \term{Task} assegnate e le relative scadenze, garantendo la tracciabilità delle scelte progettuali.
\glossentry{Vercel}{}
\term{Vercel} è una piattaforma cloud di sviluppo e hosting specializzata in applicazioni web frontend e serverless (in particolare per progetti basati su \term{Next.js}). Viene utilizzata dal \term{Team} per il deploy continuo, l'anteprima delle modifiche e la messa online della versione di produzione del software.
\glossentry{Versione}{versioni}
La \term{Versione} indica lo stato incrementale e identificabile di un \term{Documento} o di un prodotto software, gestito tramite un sistema di versionamento (come \term{Git}) per tracciarne l'evoluzione, le modifiche e la conformità alle specifiche approvate nelle diverse fasi del progetto.
\glossentry{Visual Studio Code}{VS Code, VSCode}
\term{Visual Studio Code} (VS Code) è un ambiente di sviluppo integrato (\term{Integrated Development Environment}) leggero e open source sviluppato da Microsoft, dotato di estensioni per la formattazione in \term{LaTeX}, il controllo di versione tramite \term{Git} e la scrittura di codice in \term{TypeScript}, scelto dal \term{Team} come editor principale per le attività di sviluppo.
\section{W}
\glossentry{Walkthrough}{}
Il \term{Walkthrough} è un'ispezione informale e strutturata del codice o della documentazione, condotta dai membri del \term{Team} con l'obiettivo di individuare difetti, verificare la conformità agli standard stilistici e normativi, e condividere la conoscenza tecnica prima delle fasi formali di verifica.
\glossentry{Webhook}{webhooks}
Un \term{Webhook} è un metodo di comunicazione tra applicazioni basato su eventi, che consente a un sistema di inviare una notifica HTTP automatica (solitamente in formato \term{JavaScript Object Notation}) a un URL specifico di un'altra applicazione non appena si verifica un determinato evento, facilitando l'integrazione continua e i flussi di lavoro automatizzati su piattaforme come \term{GitHub}.
